%% ============================================================
%%  Chapter 23 --- Observability with Langfuse
%%  Production-Grade Agentic AI: LangGraph + Python Frameworks
%%  Dependencies: langfuse >= 3.0.0,
%%                langchain-core >= 0.3.0,
%%                langgraph >= 0.2.0,
%%                prometheus-client >= 0.20.0,
%%                pydantic >= 2.7.0,
%%                python >= 3.11
%% ============================================================

\chapterquote{Invisible cost is not a billing problem. It is a design
problem that billing makes visible.}{anon}

The production failure this chapter prevents is undetected cost
accumulation in agentic cycles. An agent that executes thirteen nodes
across a retrieval-augmented loop returns a correct result, raises no
exception, and leaves no trace in application logs unless instrumentation
was deliberately wired before the first deployment. The economic damage
is not speculative: a single misconfigured routing edge that adds two
extra LLM calls per session can double daily token spend before any
engineer notices the invoice. The secondary failure is undifferentiated
latency---a p99 slowdown that appears in the service dashboard cannot be
attributed to a retrieval node, a long-running tool call, or a prompt
that quietly grew by eight hundred tokens after a content-policy change,
because without per-node traces there is no data to query.

The solution has three complementary layers. The first is
\emph{semantic tracing}: every LLM generation and every custom helper
function emits a Langfuse observation that records its input, output,
token counts, and prompt version, enabling Langfuse to compute
per-trace cost automatically and to answer ``which prompt version drove
this spike?''. The second is \emph{structural tracing}: LangGraph's
callback system is wired so that every node invocation produces a
Langfuse span without any instrumentation code inside node functions,
and every invocation carries the thirteen mandatory fields of the cycle
trace contract established in Chapter~1. The third is \emph{operational
metrics}: a small set of Prometheus instruments---cycle count, error
rate, HITL interrupt frequency, node latency---answer fleet-level
rate questions that a trace viewer is not designed to answer, and a
version-controlled Grafana dashboard makes those answers immediately
visible to on-call engineers.

The sections are ordered by dependency. Section~\ref{sec:obs-sdk-primitives}
establishes the Langfuse SDK primitives so that every later section can
call them with understood behaviour. Section~\ref{sec:obs-graph-tracing}
wires the LangGraph callback integration, because the handler must
exist before any node runs. Section~\ref{sec:obs-cycle-record} maps
the thirteen mandatory cycle-trace fields onto Langfuse observations,
because later sections read those fields for cost and audit purposes.
Section~\ref{sec:obs-cost-tracking} shows how token counts flow from
the LLM response into Langfuse's cost dashboard.
Section~\ref{sec:obs-prompt-management} replaces inline prompts with
versioned registry entries and links each generation to its prompt.
Section~\ref{sec:obs-memory-audit} routes the memory governance events
from Chapter~8 into the same Langfuse trace.
Sections~\ref{sec:obs-prometheus} and~\ref{sec:obs-grafana} add
the Prometheus layer and the four-panel Grafana dashboard.
Section~\ref{sec:obs-full-integration} assembles all layers into a
single runnable instrumented graph.

This chapter introduces seven terms defined in Table~\ref{tab:ch23-terms}.

\begin{table}[ht]
\centering
\caption{Observability terms introduced in Chapter~23}
\label{tab:ch23-terms}
\begin{tabularx}{\textwidth}{lX}
\toprule
\textbf{Term} & \textbf{Definition} \\
\midrule
\emph{Observation} &
  A Langfuse record of a discrete unit of work: a span, a generation,
  or an event. Observations nest to form a trace tree; the root
  observation is the trace itself. \\
\emph{Trace} &
  The root container in Langfuse that groups all observations produced
  by a single agent invocation. Identified by a \code{trace\_id} that
  can be pre-seeded from the caller. \\
\emph{Semantic tracing} &
  Tracing the intent and output of a function, usually an LLM call or
  business helper. Produced by the \code{@observe()} decorator or the
  \code{start\_as\_current\_observation()} context manager. \\
\emph{Structural tracing} &
  Tracing the lifecycle boundary of a LangGraph node. Produced by
  \code{CallbackHandler} injected via \code{RunnableConfig}; requires
  no code inside node functions. \\
\emph{Cycle trace record} &
  A typed dataclass carrying the thirteen mandatory fields that every
  node execution must record; defined in Chapter~1~\S1.9 as the
  system-wide audit contract. \\
\emph{Usage object} &
  The dict \code{\{"input": int, "output": int, "total": int,
  "unit": "TOKENS"\}} that Langfuse reads to compute generation cost
  automatically from its model-price catalogue. \\
\emph{Payload hash} &
  A 16-character hex prefix of a SHA-256 digest over the
  JSON-serialised node input or output. Stored in place of the full
  payload to bound Langfuse storage without losing identity. \\
\bottomrule
\end{tabularx}
\end{table}

\medskip
\noindent\textbf{Source layout for this chapter:}
\begin{verbatim}
src/
  observability/
    langfuse_setup.py        # §23.1 — client factory, observe patterns
    graph_runner.py          # §23.2 — CallbackHandler + RunnableConfig
    cycle_trace.py           # §23.3 — CycleTraceRecord + attachment
    cost_tracking.py         # §23.4 — usage extraction + update
    prompt_registry.py       # §23.5 — versioned prompt fetch
    memory_audit_bridge.py   # §23.6 — audit event routing
    prometheus_metrics.py    # §23.7 — Counter + Histogram instruments
  agent/
    instrumented_node.py     # §23.9 — complete instrumented node
    graph_assembly.py        # §23.9 — graph wiring + invoke helper
deploy/
  grafana/
    dashboards/
      agent_observability.json  # §23.8 — four-panel provisioned dashboard
\end{verbatim}


%% ------------------------------------------------------------
\section{Every Observation Needs an Explicit Owner}
\label{sec:obs-sdk-primitives}
%% ------------------------------------------------------------

The temptation when instrumenting an agent for the first time is to
reach for Python's \code{logging} module, add INFO lines at node entry
and exit, and treat structured tracing as a later-sprint concern. That
path fails because log lines have no native concept of parent--child
relationships, no cost-calculation engine, and no prompt-version
linkage. By the time the cost problem is undeniable, retrofitting
structured tracing into a logging-oriented codebase requires changing
every function signature and adding context propagation where none was
designed. The governing principle is that every unit of work must have
an explicit observation owner from day one---a Langfuse span or
generation object whose lifetime exactly matches the work it describes.

Langfuse's Python SDK (version~3 and later) provides three interoperable
ways to create observations, each suited to a distinct context. The
\emph{observe decorator} wraps a function and automatically records its
arguments, return value, timing, and any unhandled exception.
The \emph{context manager}, \code{langfuse.start\_as\_current\_observation()},
opens a span that becomes the active OpenTelemetry context for its
\code{with} block, automatically parenting all nested observations created
inside it; the \code{as gen} reference gives the caller a live span object
on which to call \code{update()} mid-execution. The \emph{manual
observation}, \code{langfuse.start\_observation()}, creates a span
without shifting the active context, which is the right choice for
background jobs and APScheduler tasks that run outside any
request-scoped context and must not be nested under a non-existent
parent.

All three approaches are interoperable across async boundaries. The SDK
uses OpenTelemetry context propagation internally, so a \code{@observe()}
-decorated function that calls a context-manager block inside it produces
a correctly parented hierarchy without any manual thread-local management.

The \code{Langfuse} client must be constructed once per process and shared
everywhere. Its constructor reads credentials from the environment
variables \code{LANGFUSE\_PUBLIC\_KEY}, \code{LANGFUSE\_SECRET\_KEY}, and
\code{LANGFUSE\_HOST}, so no secret appears in source code. At process
exit the SDK registers an \code{atexit} hook automatically, but services
should also call \code{langfuse.shutdown()} in the application's explicit
shutdown path so that the background ingestion buffer is flushed before
the process is replaced rather than after.

\begin{lstlisting}[language=Python,
  caption={Singleton client factory; three observation patterns},
  label={lst:obs-sdk-primitives}]
# src/observability/langfuse_setup.py
import atexit
from functools import lru_cache
from langfuse import Langfuse
from langfuse.decorators import observe

@lru_cache(maxsize=1)
def get_langfuse() -> Langfuse:
    return Langfuse()   # reads LANGFUSE_* env vars automatically

def shutdown_langfuse() -> None:
    get_langfuse().shutdown()

atexit.register(shutdown_langfuse)

# Pattern A: decorator — zero-touch function instrumentation
@observe()
def classify_intent(user_message: str, llm) -> str:
    return llm.invoke(f"Classify: {user_message}").content

# Pattern B: context manager — mid-execution metadata updates
def run_retrieval(query: str, docs: list) -> None:
    lf = get_langfuse()
    with lf.start_as_current_observation(
        name    = "retrieval",
        as_type = "span",
        input   = {"query": query},
    ) as span:
        # ... retrieval work ...
        span.update(metadata={"doc_count": len(docs)})

# Pattern C: manual observation — background / scheduler context
def emit_background_event(memory_id: str) -> None:
    lf  = get_langfuse()
    obs = lf.start_observation(name="expiry_check", as_type="event")
    # ... background work ...
    obs.end(output={"memory_id": memory_id})
\end{lstlisting}

Three decisions in Listing~\ref{lst:obs-sdk-primitives} deserve
explanation. First, \code{get\_langfuse} is wrapped in
\code{@lru\_cache(maxsize=1)}. The \code{Langfuse} constructor starts a
background ingestion thread and initialises an OpenTelemetry tracer
provider; constructing it per-request would create a new thread on every
invocation, eventually exhausting the thread pool under load. The cache
ensures exactly one client instance per process regardless of how many
modules import \code{get\_langfuse}. Second, \code{shutdown\_langfuse}
is registered with \code{atexit} in the module body, not only in a
FastAPI lifespan handler. This guarantees the flush fires in test
processes, CLI scripts, and worker processes where a lifespan handler
may never be wired. Third, \code{span.update()} is called after the
retrieval work completes rather than at construction time because only
then is \code{doc\_count} known. The SDK accumulates updates and
serialises them as a single payload when the context manager exits, so
the timing of \code{update()} within the block has no effect on what
Langfuse receives; the pattern reads naturally as ``enrich the span
with data produced during the work it describes''.

\begin{warningbox}
\code{@observe()} captures the decorated function's argument list as the
observation input. If the function receives a raw \code{AgentState}
dictionary or a large \code{AIMessage} list, the full object is
serialised and sent to Langfuse on every call. Wrap expensive arguments
in a summary or pass the \code{capture\_input=False} flag to the
decorator to avoid uploading kilobytes of state per node invocation.
\end{warningbox}

Semantic tracing covers individual functions. The next problem is
tracing the structural shape of the LangGraph execution itself---the
sequence of nodes, their durations, and their parent--child
relationships---without writing any instrumentation code inside node
bodies.


%% ------------------------------------------------------------
\section{Graph Structure Must Be Traced Outside Node Bodies}
\label{sec:obs-graph-tracing}
%% ------------------------------------------------------------

The naive approach to LangGraph observability is to add a Langfuse call
at the start and end of every node function. This couples observability
to business logic in a way that is easy to introduce and nearly
impossible to remove cleanly: every new node must receive the tracing
boilerplate, every tracing-strategy change requires editing every node,
and any node that is omitted from the refactor produces silent trace
gaps. The governing principle is that structural tracing belongs at the
invocation boundary---between LangGraph and the node function---not
inside the function itself.

LangGraph threads a \code{RunnableConfig} object through every node
invocation. The config carries a \code{callbacks} list whose handlers
are notified at the start and end of each runnable's execution.
Langfuse provides \code{CallbackHandler}---importable from
\code{langfuse.langchain}---that implements LangChain's callback
protocol and emits one Langfuse span per node invocation,
automatically nesting child observations (embedded LLM calls, tool
calls) under the correct parent span. The handler is passed to
\code{graph.invoke()} or \code{graph.stream()} via the \code{callbacks}
field of \code{RunnableConfig}; it is not a constructor argument to
\code{StateGraph} and it is not attached to individual nodes.

To bind session identity to every span in the invocation, call
\code{langfuse.propagate\_attributes()} before the graph runs.
Langfuse writes these values into the active OpenTelemetry context, and
every observation created downstream---including those emitted by
\code{CallbackHandler} for each node---inherits them automatically.

\begin{warningbox}
Construct a new \code{CallbackHandler} instance per graph invocation,
never per process. A shared handler instance accumulates open-span
state; when two concurrent requests share the same instance, their spans
interleave inside a single Langfuse trace, producing a trace tree that
is structurally invalid and unreadable. The handler is lightweight: its
constructor performs no I/O, so creating one per request has no
meaningful overhead.
\end{warningbox}

\begin{lstlisting}[language=Python,
  caption={Per-invocation \code{CallbackHandler} and attribute
  propagation via \code{RunnableConfig}},
  label={lst:obs-graph-tracing}]
# src/observability/graph_runner.py
from langchain_core.runnables import RunnableConfig
from langfuse.langchain import CallbackHandler
from src.observability.langfuse_setup import get_langfuse

def run_graph_with_tracing(
    graph,
    state:      dict,
    user_id:    str,
    session_id: str,
) -> dict:
    lf      = get_langfuse()
    handler = CallbackHandler()      # fresh instance per call
    lf.propagate_attributes(
        user_id    = user_id,
        session_id = session_id,
    )
    config = RunnableConfig(
        callbacks = [handler],
        metadata  = {
            "user_id":    user_id,
            "session_id": session_id,
        },
    )
    try:
        return graph.invoke(state, config=config)
    finally:
        lf.shutdown()
\end{lstlisting}

Two decisions in Listing~\ref{lst:obs-graph-tracing} deserve explanation.
First, \code{lf.propagate\_attributes()} is called \emph{before}
\code{graph.invoke()}, not after. The method writes values into the
active OpenTelemetry context, which is inherited by every observation
created during the downstream call. Calling it after \code{graph.invoke()}
would inject attributes into a context that is already fully unwound,
leaving every node span with a missing \code{user\_id} and
\code{session\_id}. Second, \code{metadata} is populated in
\code{RunnableConfig} even though \code{propagate\_attributes()} already
handles the Langfuse path. \code{RunnableConfig.metadata} is consumed
by LangGraph's Studio introspection, LangSmith tracing, and any node
that calls \code{config.get("metadata", \{\})} to read request-scoped
values; it must be present for full portability across the LangChain
tooling stack.

Graph nodes now emit structural spans automatically. The remaining
problem is that those spans carry whatever metadata an individual developer
happened to add---which, without a shared contract, produces an
inconsistent trace corpus where cost queries and latency queries fail on
nodes that simply forgot a field.


%% ------------------------------------------------------------
\section{Every Node Execution Must Carry a Standard Contract}
\label{sec:obs-cycle-record}
%% ------------------------------------------------------------

When engineers instrument nodes independently, each adds the fields they
consider relevant, and the result is a Langfuse instance where some
spans carry latency, some carry a model name, and none carry all
thirteen fields needed to answer the question ``which node, on which
cycle, for which user, caused this cost spike?''. A style guide is not
enforced at runtime and degrades under deadline pressure. The governing
principle is that the span metadata contract is a typed dataclass
assembled in a single factory function, making compliance automatic
rather than cultural.

Chapter~1~\S1.9 defined the thirteen-field \code{CycleTraceRecord}
as the book's system-wide audit contract. The first four fields are
identity and scoping fields that belong on the trace root:
\code{trace\_id}, \code{cycle\_id}, \code{session\_id}, and
\code{user\_id}. The next four are node-execution fields that belong on
the node span: \code{node\_name}, \code{node\_input\_hash},
\code{node\_output\_hash}, and \code{started\_at}. The final five
capture generation performance: \code{ended\_at}, \code{latency\_ms},
\code{llm\_model}, \code{input\_tokens}, and \code{output\_tokens}.

Payload hashing deserves its own design explanation. Storing the full
node input and output in Langfuse metadata would bloat every span with
potentially thousands of tokens of JSON, making the Langfuse UI slow
and storage costs significant. \code{node\_input\_hash} and
\code{node\_output\_hash} are 16-character hex prefixes of SHA-256
digests over the JSON-serialised payload. This is enough to verify
identity---two identical inputs produce the same hash, enabling
deduplication queries---without storing the payload itself. When the
full payload is needed for debugging, it is retrievable from the
LangGraph checkpoint store via \code{AsyncPostgresSaver.aget()}, which
was established in Chapter~22~\S22.1.

\begin{lstlisting}[language=Python,
  caption={\code{CycleTraceRecord} dataclass and span attachment helper},
  label={lst:obs-cycle-record}]
# src/observability/cycle_trace.py
import hashlib, json
from dataclasses import asdict, dataclass
from datetime import datetime, timezone
from typing import Any

@dataclass
class CycleTraceRecord:
    trace_id:          str
    cycle_id:          str
    session_id:        str
    user_id:           str
    node_name:         str
    node_input_hash:   str
    node_output_hash:  str
    started_at:        datetime
    ended_at:          datetime
    latency_ms:        float
    llm_model:         str
    input_tokens:      int
    output_tokens:     int

def _hash16(obj: Any) -> str:
    raw = json.dumps(obj, sort_keys=True, default=str)
    return hashlib.sha256(raw.encode()).hexdigest()[:16]

def build_cycle_trace(
    state:     dict,
    node_name: str,
    start:     datetime,
    end:       datetime,
    model:     str,
    usage:     dict[str, int],
) -> CycleTraceRecord:
    meta = state.get("meta", {})
    return CycleTraceRecord(
        trace_id         = meta.get("trace_id", ""),
        cycle_id         = meta.get("cycle_id", ""),
        session_id       = meta.get("session_id", ""),
        user_id          = meta.get("user_id", ""),
        node_name        = node_name,
        node_input_hash  = _hash16(state.get("messages", [])),
        node_output_hash = _hash16(state.get("output", "")),
        started_at       = start,
        ended_at         = end,
        latency_ms       = (end - start).total_seconds() * 1000,
        llm_model        = model,
        input_tokens     = usage.get("input_tokens", 0),
        output_tokens    = usage.get("output_tokens", 0),
    )
\end{lstlisting}

Two decisions in Listing~\ref{lst:obs-cycle-record} deserve explanation.
First, \code{\_hash16} passes \code{default=str} to \code{json.dumps}.
Real agent state dictionaries contain \code{datetime} objects, Pydantic
models, and LangChain \code{AIMessage} instances---none of which the
default JSON encoder handles. Passing \code{default=str} serialises
them as their string representations, which is lossy for round-trip
fidelity but entirely correct for the hash's purpose: verifying
identity, not reconstructing the object. Second, \code{started\_at} and
\code{ended\_at} are captured by the caller using
\code{datetime.now(timezone.utc)} at the actual moment of node entry
and exit, not from Langfuse's server-side timestamps. Langfuse's
timestamps record when the span payload arrived at the ingestion API,
which may be seconds or minutes after the node ran if the background
flush queue is congested under load.

The trace record captures what happened and when. The remaining problem
is making Langfuse aware of what each generation \emph{cost}, so the
cost dashboard activates automatically rather than requiring a separate
accounting pipeline.


%% ------------------------------------------------------------
\section{Cost Must Flow Automatically from the Model Response}
\label{sec:obs-cost-tracking}
%% ------------------------------------------------------------

The naive approach to cost tracking is a side-channel database table
that stores token counts per request, joined with a price list managed
as a configuration file. This architecture requires two synchronised
updates every time a model changes its pricing---the price file and the
accounting query---and the updates can arrive out of order, producing
a window where costs are under- or over-reported. Langfuse eliminates
the side channel: if a generation observation carries a \code{usage}
object in the canonical format and a \code{model} string that matches
an entry in Langfuse's model-price catalogue, cost is computed and
aggregated automatically with no additional application code. The
governing principle is that token counts belong on the generation
observation in the exact format the platform expects, not in a
secondary store.

Langfuse's generation span exposes an \code{update()} method that
accepts a \code{usage} parameter. That parameter must be a dict with
the keys \code{"input"}, \code{"output"}, \code{"total"}, and
\code{"unit"}. The value of \code{"unit"} must be the string literal
\code{"TOKENS"}; any other value is treated as an opaque label and
bypasses the cost-computation engine entirely---there is no error, just
a silent gap in the dashboard. The \code{model} parameter must match a
name in Langfuse's public model catalogue, or a custom entry registered
via \code{langfuse.create\_model()} for fine-tuned or privately hosted
models.

LangChain and LiteLLM expose token counts on different response
attributes. LangChain's \code{AIMessage} carries a \code{usage\_metadata}
dict with keys \code{"input\_tokens"} and \code{"output\_tokens"}.
LiteLLM's \code{ModelResponse} carries a \code{usage} object with
attributes \code{prompt\_tokens} and \code{completion\_tokens}. A
single extraction function normalises both shapes so that no node
function needs to branch on the LLM provider.

\begin{lstlisting}[language=Python,
  caption={Usage extraction from two LLM response shapes; canonical
  usage object construction; generation span update},
  label={lst:obs-cost-tracking}]
# src/observability/cost_tracking.py
from typing import Any

def extract_usage(response: Any) -> dict[str, int]:
    # LangChain AIMessage.usage_metadata
    if hasattr(response, "usage_metadata") and response.usage_metadata:
        m = response.usage_metadata
        return {
            "input_tokens":  m.get("input_tokens", 0) or 0,
            "output_tokens": m.get("output_tokens", 0) or 0,
        }
    # LiteLLM ModelResponse.usage
    if hasattr(response, "usage") and response.usage:
        u = response.usage
        return {
            "input_tokens":  getattr(u, "prompt_tokens", 0) or 0,
            "output_tokens": getattr(u, "completion_tokens", 0) or 0,
        }
    return {"input_tokens": 0, "output_tokens": 0}

def build_usage_object(raw: dict[str, int]) -> dict:
    total = raw["input_tokens"] + raw["output_tokens"]
    return {
        "input":  raw["input_tokens"],
        "output": raw["output_tokens"],
        "total":  total,
        "unit":   "TOKENS",
    }

def update_generation_span(
    span, model: str, response: Any
) -> dict[str, int]:
    raw = extract_usage(response)
    span.update(model=model, usage=build_usage_object(raw))
    return raw
\end{lstlisting}

Three decisions in Listing~\ref{lst:obs-cost-tracking} deserve
explanation. First, both \code{getattr} calls and both \code{dict.get}
calls use \code{or 0} after the fallback default. LiteLLM's
\code{usage} fields and LangChain's \code{usage\_metadata} fields are
optional; when the upstream provider does not return counts (as with
some streaming responses where the final chunk has not yet arrived)
the fields may be \code{None} rather than absent, and
\code{getattr(u, "prompt\_tokens", 0)} on a \code{None}-valued
attribute returns \code{None}, not \code{0}. The \code{or 0}
guard ensures that downstream arithmetic never receives \code{None}.
Second, \code{update\_generation\_span} returns the raw token dict.
This allows the same call to satisfy both the Langfuse span and the
\code{build\_cycle\_trace} factory from Section~\ref{sec:obs-cycle-record}
without calling \code{extract\_usage} twice. Third, \code{build\_usage\_object}
is a pure function separate from the span-update logic. This makes it
testable in isolation: a unit test can assert that a given raw dict
produces the correct \code{"TOKENS"} object without requiring a live
Langfuse client.

Token counts reach Langfuse automatically. The remaining problem is
that raw token counts do not identify \emph{which prompt version}
drove a cost change, because a prompt edit that doubles average output
length is indistinguishable from a routing change that doubled call
frequency unless the prompt's identity is recorded on the generation
observation.


%% ------------------------------------------------------------
\section{Prompts Must Be Versioned Before They Reach the Model}
\label{sec:obs-prompt-management}
%% ------------------------------------------------------------

A prompt stored as an inline Python f-string has no version, no
deployment history, and no linkage to the generation that used it.
When output quality degrades after an informal edit, the only
diagnostic tool is \code{git blame}---which works only if the change
was committed, cannot identify which production generations used the
bad version, and offers no rollback that does not require a deployment.
The governing principle is that prompts are configuration, not code:
they must be stored in a registry, fetched at runtime by label, and
attached to every generation observation so Langfuse can aggregate
performance by prompt version without any additional application logic.

Langfuse's prompt registry stores named, versioned templates.
\code{langfuse.get\_prompt(name, label="production")} fetches the
template currently tagged as production, caches it locally with a
configurable TTL (default 60 seconds), and returns a
\code{TextPromptClient} or \code{ChatPromptClient} object.
\code{prompt.compile(**variables)} substitutes placeholders and returns
the final string. \code{prompt.get\_langchain\_prompt()} returns a
LangChain \code{ChatPromptTemplate} whose variable syntax is normalised
from Langfuse's single-brace format to LangChain's double-brace format.

\begin{warningbox}
Langfuse template variables use single braces: \code{\{variable\}}.
LangChain's \code{ChatPromptTemplate} uses double braces for its own
variables. Always call \code{prompt.get\_langchain\_prompt()} when the
result will be used with a LangChain chain; using \code{compile()} and
passing the raw string will leave un-substituted placeholder text in the
message that reaches the model---silently and without any error from
LangChain.
\end{warningbox}

Attaching the prompt object to the generation is the step most
developers omit. The attachment is a single \code{update()} keyword
argument---\code{span.update(prompt=prompt\_client\_object)}---and once
present, the Langfuse UI can filter all generations by prompt name and
version, plot latency and cost per version side by side, and display a
diff between any two versions. Without the attachment the prompt data
exists in Langfuse's registry but the per-generation linkage is absent,
and the cost dashboard shows only ``no associated prompt'' for every
affected generation.

\begin{lstlisting}[language=Python,
  caption={Production-labelled prompt fetch with 5-minute cache;
  LangChain-compatible variant; span attachment helper},
  label={lst:obs-prompt-management}]
# src/observability/prompt_registry.py
from langfuse import Langfuse
from src.observability.langfuse_setup import get_langfuse

def get_production_prompt(
    name:      str,
    variables: dict,
):
    lf     = get_langfuse()
    prompt = lf.get_prompt(
        name,
        label             = "production",
        cache_ttl_seconds = 300,
    )
    compiled: str = prompt.compile(**variables)
    return compiled, prompt   # caller uses prompt for span attachment

def get_langchain_prompt(name: str):
    lf     = get_langfuse()
    prompt = lf.get_prompt(
        name,
        label             = "production",
        cache_ttl_seconds = 300,
    )
    return prompt.get_langchain_prompt(), prompt

def attach_prompt(span, prompt) -> None:
    span.update(prompt=prompt)
\end{lstlisting}

Two decisions in Listing~\ref{lst:obs-prompt-management} deserve
explanation. First, \code{cache\_ttl\_seconds=300} overrides the
60-second default. The default is calibrated for development, where
prompt edits should be visible quickly. In production, a 60-second
window means an incident-response prompt rollback may not reach all
running instances for up to a minute; 300 seconds provides a
predictable, documentable window without requiring per-instance
cache-invalidation signals. Second, both \code{get\_production\_prompt}
and \code{get\_langchain\_prompt} return a \code{(compiled, prompt)}
pair rather than only the compiled string. A function that returns only
the string forces the caller to issue a second \code{get\_prompt()} call
to obtain the object needed for \code{attach\_prompt}, producing an
extra cache lookup or an extra API call during every node execution.

Prompt versions are now linked to every generation. The remaining problem
is a different class of observation: the memory governance events from
Chapter~8, which must appear in the same Langfuse trace as the agent
session that triggered them rather than as disconnected top-level traces.


%% ------------------------------------------------------------
\section{Audit Events Must Flow Into the Active Trace}
\label{sec:obs-memory-audit}
%% ------------------------------------------------------------

When memory write and retrieval operations are emitted as top-level
Langfuse traces rather than as observations nested under the current
agent session, the Langfuse UI shows a fragmented picture: session
traces with no memory context, and memory traces with no session
context. Correlating them requires exporting both datasets and joining
on \code{session\_id}---a manual operation that should never be
necessary. The governing principle is that memory audit events are
observations in the same trace as the agent invocation that caused them,
not records in a parallel trace.

Chapter~8~\S8.12 defined three audit dataclasses and their Langfuse
emitter functions: \code{WriteAuditRecord} carries the outcome of a
governed memory write; \code{RetrievalAuditRecord} carries the outcome
of a narrowed memory retrieval; \code{ExpiryAuditRecord} carries the
outcome of a scheduled deletion. Those classes exist in
\code{src.agent.memory.audit} and must not be redefined here.
Redefining them would create two schemas for the same concept that can
drift apart silently when Chapter~8's module is updated.

The bridge module in this chapter adds only the context-threading logic
that Chapter~8 did not yet need: how to obtain the active Langfuse
trace ID from the OpenTelemetry context and pass it to the emitter
functions so they nest their observations under the right parent.
\code{langfuse.get\_current\_trace\_id()} returns the trace ID of the
observation currently active in the calling thread's context; inside a
\code{graph.invoke()} call where \code{propagate\_attributes()} was
called beforehand, this returns the session's root trace ID.

\begin{notebox}
\code{ExpiryAuditRecord} events are emitted by the APScheduler background
job defined in Chapter~8~\S8.8, which runs in a separate thread outside
any request-scoped OpenTelemetry context. Calling
\code{get\_current\_trace\_id()} in that thread returns an empty string.
These events must be emitted as standalone observations with no parent
trace; passing an empty string as \code{trace\_id} to the Langfuse API
creates an orphaned observation that cannot be located through normal
session queries in some Langfuse versions.
\end{notebox}

\begin{lstlisting}[language=Python,
  caption={Audit bridge: trace-ID injection for session context;
  background-safe standalone emission for expiry events},
  label={lst:obs-memory-audit}]
# src/observability/memory_audit_bridge.py
from src.agent.memory.audit import (
    ExpiryAuditRecord,
    RetrievalAuditRecord,
    WriteAuditRecord,
    emit_expiry_audit,
    emit_retrieval_audit,
    emit_write_audit,
)
from src.observability.langfuse_setup import get_langfuse

def emit_write_event(record: WriteAuditRecord) -> None:
    lf       = get_langfuse()
    trace_id = lf.get_current_trace_id()
    emit_write_audit(record=record, lf=lf, trace_id=trace_id)

def emit_retrieval_event(record: RetrievalAuditRecord) -> None:
    lf       = get_langfuse()
    trace_id = lf.get_current_trace_id()
    emit_retrieval_audit(record=record, lf=lf, trace_id=trace_id)

def emit_expiry_event(record: ExpiryAuditRecord) -> None:
    lf = get_langfuse()
    # No trace_id: APScheduler thread has no active span.
    emit_expiry_audit(record=record, lf=lf)
\end{lstlisting}

Two decisions in Listing~\ref{lst:obs-memory-audit} deserve explanation.
First, the three functions call \code{get\_langfuse()} independently
rather than accepting a \code{Langfuse} argument. This would normally
violate explicit dependency injection, but the \code{@lru\_cache}
singleton in \code{langfuse\_setup} guarantees that all three calls
return the same instance, so there is no resource overhead. More
importantly, passing the client as an argument would expose the
Langfuse dependency to every memory node function, coupling the memory
subsystem to the observability subsystem. The bridge module is the
only place that knows about both; calling the singleton keeps that
coupling localised there. Second, \code{emit\_expiry\_event} omits the
\code{get\_current\_trace\_id()} call entirely rather than calling it
and discarding an empty string. In the APScheduler thread context,
\code{get\_current\_trace\_id()} may return an empty string
\emph{or} the stale trace ID of the last request that happened to run
on the same OS thread before the scheduler task was scheduled---a
non-deterministic value. Omitting the call prevents accidentally
parenting an expiry event under a completely unrelated session trace.

Session events and memory events now appear in one trace. The remaining
problem is the class of question that a trace viewer is not designed to
answer: what is the error rate of the \code{tool\_call\_node} right now,
across all sessions, over the last five minutes?


%% ------------------------------------------------------------
\section{Operational Questions Require Their Own Instruments}
\label{sec:obs-prometheus}
%% ------------------------------------------------------------

Langfuse is a trace viewer and cost analyser, not a time-series
database. Its query interface is not built to answer ``what is the
p99 latency for \code{retrieval\_node} between 14:00 and 15:00 across
all users?'' in sub-second latency. Routing that question to Prometheus
---a purpose-built time-series database with PromQL designed for rate
and quantile calculations---keeps each system operating within its
design envelope. The governing principle is that trace observations
answer ``what happened in this session'' while Prometheus metrics answer
``what is the fleet doing right now''.

This chapter uses \code{prometheus\_client.Counter} and
\code{prometheus\_client.Histogram} only, matching the constraint of
the official Python client library. No custom aggregation class, no
\code{Gauge} for values that reset on process restart (those produce
misleading spikes after a redeploy), and no \code{user\_id} label
(label cardinality grows with the user population, eventually making
Prometheus scrape latency unacceptable). The label set for all four
instruments is exactly \code{node\_name}, whose cardinality is bounded
by the number of nodes in the graph.

Four instruments cover the operational questions that matter.
\code{agent\_cycle\_total} counts node executions and supports
per-node throughput queries. \code{agent\_cycle\_error\_total} counts
unhandled exceptions and, divided by \code{agent\_cycle\_total}, gives
the per-node error rate. \code{agent\_node\_latency\_seconds} is a
histogram calibrated to LLM operation latencies, enabling p99
computation via \code{histogram\_quantile}. \code{agent\_hitl\_interrupt\_total}
counts calls to LangGraph's \code{interrupt()} primitive---introduced
in Chapter~21 as the mechanism by which a node pauses execution for
human review; an elevated rate signals that the review queue is growing
faster than it is being processed.

\begin{lstlisting}[language=Python,
  caption={Four Prometheus instruments and the \code{instrument\_node}
  decorator},
  label={lst:obs-prometheus}]
# src/observability/prometheus_metrics.py
import time
from functools import wraps
from typing import Callable, TypeVar
from prometheus_client import Counter, Histogram, start_http_server

CYCLE_COUNT = Counter(
    "agent_cycle_total",
    "Total node executions",
    ["node_name"],
)
ERROR_COUNT = Counter(
    "agent_cycle_error_total",
    "Unhandled exceptions per node",
    ["node_name"],
)
NODE_LATENCY = Histogram(
    "agent_node_latency_seconds",
    "Node wall-clock latency in seconds",
    ["node_name"],
    buckets=[0.01, 0.05, 0.1, 0.25, 0.5, 1.0, 2.5, 5.0, 10.0],
)
HITL_COUNT = Counter(
    "agent_hitl_interrupt_total",
    "LangGraph interrupt() calls per node",
    ["node_name"],
)

F = TypeVar("F", bound=Callable)

def instrument_node(node_name: str) -> Callable[[F], F]:
    def decorator(fn: F) -> F:
        @wraps(fn)
        def wrapper(*args, **kwargs):
            CYCLE_COUNT.labels(node_name).inc()
            t0 = time.perf_counter()
            try:
                return fn(*args, **kwargs)
            except Exception:
                ERROR_COUNT.labels(node_name).inc()
                raise
            finally:
                NODE_LATENCY.labels(node_name).observe(
                    time.perf_counter() - t0
                )
        return wrapper  # type: ignore[return-value]
    return decorator

def start_metrics_server(port: int = 8001) -> None:
    start_http_server(port)
\end{lstlisting}

Three decisions in Listing~\ref{lst:obs-prometheus} deserve explanation.
First, the \code{NODE\_LATENCY} histogram's highest finite bucket is
10 seconds. LLM calls backed by large models can take 5--8 seconds
under load; a histogram with a highest finite bucket of 2.5 seconds
would place all slow calls in the \code{+Inf} overflow bucket, making
\code{histogram\_quantile(0.99,~...)} return \code{+Inf} rather than a
real latency value and rendering the p99 panel useless during the exact
incidents it is meant to diagnose. Second, \code{@wraps(fn)} preserves
the wrapped function's \code{\_\_name\_\_} and \code{\_\_doc\_\_}.
LangGraph introspects the function name to produce node labels in its
Studio UI and in error messages; without \code{@wraps}, every
instrumented node appears as ``wrapper'' in those tools. Third,
\code{ERROR\_COUNT} is incremented inside the \code{except} block
before a bare \code{raise}. The increment occurs at the exact moment
an exception escapes the node, and the original exception propagates
unmodified so LangGraph's own error handling and retry logic operates
on the real exception type.

\begin{tipbox}
Register \code{start\_metrics\_server} in the application's startup
event handler---FastAPI's \code{lifespan} context or an
\code{@app.on\_event("startup")} hook---rather than at module import
time. Importing \code{prometheus\_metrics} during a test run would
start an HTTP server on port 8001, causing port-conflict failures on CI
runners that execute test workers in parallel.
\end{tipbox}

Prometheus now captures rate, error, and latency data. The remaining
problem is presenting that data in a reproducible form that supports
real-time monitoring and threshold-based alerting without requiring
on-call engineers to write PromQL from scratch when they open the
dashboard during an incident.


%% ------------------------------------------------------------
\section{Dashboards Are Code, Not Click-Ops}
\label{sec:obs-grafana}
%% ------------------------------------------------------------

A Grafana dashboard built through the UI is not reproducible. It lives
in the Grafana database, is not version-controlled, cannot be reviewed
in a pull request, and is silently lost when the Grafana instance is
replaced without a volume backup. The governing principle is that the
dashboard is a JSON file committed to the repository and provisioned
via Grafana's file-based provisioning system---mounted as a Kubernetes
ConfigMap or a Docker Compose bind mount---so it is deployed with the
application, reviewed like code, and recoverable from source control.

The dashboard declares a single Prometheus datasource identified by UID.
Grafana's provisioning system resolves datasources by UID, not by
name; hardcoding the name ``Prometheus'' works only if the provisioned
datasource carries that exact name in every deployment environment, which
is not guaranteed. The UID is stable across Grafana upgrades and
datasource renames and is the correct anchor for cross-environment
reproducibility.

Four panels map directly onto the four instruments from
Section~\ref{sec:obs-prometheus}. Each description below specifies
the panel title, the PromQL query, the visualisation type, and the
alert condition precisely enough to reproduce the panel from the JSON
specification without a screenshot.

\textbf{Panel~1 --- Cycle Rate by Node.} Title: ``Cycle Rate (req/s)''.
Query: \code{rate(agent\_cycle\_total\{job="agent"\}[5m])} with
\code{by~(node\_name)} grouping. Visualisation: time-series, one
series per node, 30-second dashboard refresh. No alert threshold; this
panel provides the denominator for error-rate calculations and a
capacity-planning baseline.

\textbf{Panel~2 --- Per-Node Error Rate.} Title: ``Error Rate
(errors/cycle)''. Query:
\begin{verbatim}
  rate(agent_cycle_error_total{job="agent"}[5m])
/ rate(agent_cycle_total{job="agent"}[5m])
\end{verbatim}
with \code{by~(node\_name)} grouping. Visualisation: time-series.
Alert condition: value \code{> 0.05} (5\,\% error rate) for two
consecutive 60-second evaluation windows, routing to PagerDuty.

\textbf{Panel~3 --- HITL Interrupt Frequency.} Title: ``HITL
Interrupts/min''. Query:
\code{rate(agent\_hitl\_interrupt\_total\{job="agent"\}[1m]) * 60}
with \code{by~(node\_name)} grouping. Visualisation: time-series.
Alert condition: value \code{> 10} per minute sustained for five
consecutive evaluation windows, routing to Slack. An elevated rate is
a scaling signal, not a correctness signal: it means the human review
queue is growing faster than reviewers can drain it.

\textbf{Panel~4 --- p99 Node Latency.} Title: ``p99 Latency (s)''.
Query:
\begin{verbatim}
  histogram_quantile(
    0.99,
    rate(agent_node_latency_seconds_bucket{job="agent"}[5m])
  )
\end{verbatim}
with \code{by~(node\_name)} grouping. Visualisation: time-series.
Alert condition: value \code{> 2.5} seconds for three consecutive
60-second evaluation windows, routing to PagerDuty.

\begin{notebox}
The \code{job="agent"} label selector in all four queries assumes that
the Prometheus scrape configuration assigns \code{job: agent} to the
scrape target for the agent process. Without the selector, queries
aggregate across every scraped service on the same Prometheus instance,
producing meaningless combined metrics when the deployment includes
other instrumented services. Chapter~26 provides the complete scrape
configuration and the Kubernetes ConfigMap that provisions this
dashboard.
\end{notebox}

Operational visibility is now covered by both Langfuse traces and
Prometheus dashboards. The final problem is assembling all layers into
a single agent whose instrumentation composes without coupling any
layer's concerns into the node functions it observes.


%% ------------------------------------------------------------
\section{Instrumentation Layers Must Compose Without Coupling}
\label{sec:obs-full-integration}
%% ------------------------------------------------------------

The risk of assembling multiple instrumentation layers is that each
new layer adds code to every node function until the instrumentation
surface exceeds the business logic. A node that opens a Langfuse
context manager, fetches a versioned prompt, calls an LLM, extracts
token counts, builds a cycle record, attaches a usage object, and
increments a Prometheus counter is executing six cross-cutting concerns
before any domain work begins. The governing principle is that
instrumentation must compose from the outside in: the Prometheus
decorator covers the node boundary, the callback handler covers the
graph boundary, and the context manager covers the LLM generation
boundary---each layer at its own level of abstraction, none reaching
into the others.

The composition works as follows. \code{@instrument\_node} from
Section~\ref{sec:obs-prometheus} wraps the node function at the
outermost boundary, recording the full node wall-clock latency that
LangGraph observes. Inside the node, an
\code{lf.start\_as\_current\_observation(as\_type="generation")}
block wraps only the LLM call, so the generation observation captures
model latency alone. The prompt is fetched from the registry
(Section~\ref{sec:obs-prompt-management}) and attached to the
generation. Token counts are extracted via \code{update\_generation\_span}
(Section~\ref{sec:obs-cost-tracking}) and passed to
\code{build\_cycle\_trace} (Section~\ref{sec:obs-cycle-record}).
The \code{CallbackHandler} from Section~\ref{sec:obs-graph-tracing}
operates entirely outside the node and requires no code changes inside
it. None of these concerns requires another to be present.

\begin{lstlisting}[language=Python,
  caption={A fully instrumented LangGraph node; all layers applied
  without coupling},
  label={lst:obs-instrumented-node}]
# src/agent/instrumented_node.py
from datetime import datetime, timezone
from langfuse import Langfuse
from src.observability.cost_tracking import update_generation_span
from src.observability.cycle_trace import build_cycle_trace
from src.observability.langfuse_setup import get_langfuse
from src.observability.prometheus_metrics import instrument_node
from src.observability.prompt_registry import (
    attach_prompt, get_production_prompt,
)

@instrument_node("llm_node")
def llm_node(state: dict) -> dict:
    lf    = get_langfuse()
    start = datetime.now(timezone.utc)
    model = state["meta"]["llm_model"]
    compiled, prompt = get_production_prompt(
        "agent_system_prompt",
        variables={"user_id": state["meta"].get("user_id", "")},
    )
    with lf.start_as_current_observation(
        name    = "llm_generation",
        as_type = "generation",
        input   = {"prompt_preview": compiled[:200]},
    ) as gen:
        response = state["llm"].invoke(compiled)
        raw      = update_generation_span(gen, model, response)
        attach_prompt(gen, prompt)
    end    = datetime.now(timezone.utc)
    record = build_cycle_trace(
        state=state, node_name="llm_node",
        start=start, end=end,
        model=model, usage=raw,
    )
    return {
        **state,
        "output": response.content,
        "meta":   {**state["meta"], "last_cycle_trace": record},
    }
\end{lstlisting}

Three decisions in Listing~\ref{lst:obs-instrumented-node} deserve
explanation. First, the \code{with~lf.start\_as\_current\_observation(...)}
block wraps only the \code{state["llm"].invoke()} call, not the
prompt fetch or cycle-record construction. The generation observation's
latency field should reflect model response time; inflating it with
prompt-registry cache lookups and JSON hashing would overstate model
latency and understate infrastructure overhead, distorting the cost
dashboard's per-generation analysis. Second, \code{attach\_prompt} is
called \emph{inside} the context-manager block, after \code{response}
is available, but this is equivalent to calling it at construction time
because the SDK accumulates \code{update()} calls and serialises them
as one payload when the block exits. The placement after the LLM call
is intentional: it reads as ``the prompt object is part of the
generation record''---which semantically it is---rather than as a
setup step that must precede the model call. Third, \code{@instrument\_node}
is the outermost decorator, wrapping the full node including prompt
fetch and record construction. This ensures the Prometheus histogram
measures total node wall-clock time as seen by LangGraph, which is the
latency that governs SLA calculations, not only the model call duration.

\begin{lstlisting}[language=Python,
  caption={Graph assembly: all instrumentation layers wired;
  per-invocation \code{invoke} helper},
  label={lst:obs-graph-assembly}]
# src/agent/graph_assembly.py
import atexit
from dataclasses import dataclass, field
from langchain_core.runnables import RunnableConfig
from langfuse.langchain import CallbackHandler
from langgraph.graph import END, START, StateGraph
from langgraph.graph.state import CompiledStateGraph
from src.agent.instrumented_node import llm_node
from src.observability.langfuse_setup import get_langfuse
from src.observability.prometheus_metrics import start_metrics_server

@dataclass
class ObservabilityConfig:
    prom_port: int = 8001

def build_instrumented_graph(
    obs: ObservabilityConfig,
) -> CompiledStateGraph:
    start_metrics_server(obs.prom_port)
    builder = StateGraph(dict)
    builder.add_node("llm_node", llm_node)
    builder.add_edge(START, "llm_node")
    builder.add_edge("llm_node", END)
    return builder.compile()

def invoke(
    graph:      CompiledStateGraph,
    state:      dict,
    user_id:    str,
    session_id: str,
) -> dict:
    lf      = get_langfuse()
    handler = CallbackHandler()
    lf.propagate_attributes(user_id=user_id, session_id=session_id)
    config = RunnableConfig(
        callbacks = [handler],
        metadata  = {"user_id": user_id, "session_id": session_id},
    )
    try:
        return graph.invoke(state, config=config)
    finally:
        lf.shutdown()
\end{lstlisting}

Two decisions in Listing~\ref{lst:obs-graph-assembly} deserve
explanation. First, \code{ObservabilityConfig} is a plain
\code{@dataclass} with a single \code{prom\_port} field. The Langfuse
credentials are consumed from environment variables inside
\code{get\_langfuse()}, so the config object has no secrets to
validate and no Pydantic model is warranted; adding one would introduce
validation overhead for a one-field object that cannot fail in any
interesting way. Second, \code{lf.shutdown()} is called in the
\code{finally} block of \code{invoke()} rather than only in the
\code{atexit} hook. In long-running API servers, \code{atexit} fires
once at process termination---potentially hours after the last
request---leaving a window where traces buffered in memory are not yet
visible in Langfuse. Calling \code{shutdown()} in \code{finally}
flushes the buffer synchronously after each request, ensuring the
trace appears before the HTTP response is returned. For
high-throughput async services where per-request flushing is too
expensive, the \code{atexit} hook alone is sufficient because the
background ingestion thread processes the queue continuously and
\code{shutdown()} in that context would block the event loop.

Chapter~24 builds the evaluation harness that imports
\code{build\_instrumented\_graph}, replays saved Langfuse traces as
test fixtures, and asserts that quality metrics do not regress across
prompt version changes---relying on the prompt linkage established in
Section~\ref{sec:obs-prompt-management} to identify exactly which
prompt version each recorded generation used.
